\documentclass[11pt]{article}
\usepackage[margin=1in]{geometry}
\usepackage{amsmath,amssymb,amsthm}
\usepackage[hidelinks]{hyperref}

\newtheorem*{theorem}{Theorem}

\newcommand{\R}{\mathbb{R}}
\newcommand{\Lcal}{\mathcal{L}}

\title{Literature Status for the Main Conjecture of arXiv:1606.06913}
\author{Automated mathematics run \texttt{fa\_banach\_001}}
\date{11 August 2026}

\begin{document}
\maketitle

\section*{Status}

\textbf{Result type:} literature already answered.

\textbf{Source:} Daniel Dadush and Oded Regev,
\emph{Towards Strong Reverse Minkowski-type Inequalities for Lattices},
arXiv:1606.06913.

\textbf{Answer:}\quad Oded Regev and Noah Stephens-Davidowitz,\\
\emph{A Reverse Minkowski Theorem}, arXiv:1611.05979.

\section*{Source conjecture}

Conjecture 3.1 (``Main conjecture'') of the source, on PDF page 14, defines
$C_\eta(n)$ to be the least constant for which every full-rank lattice
$\Lcal\subset\R^n$ satisfies
\begin{equation}
 \eta(\Lcal)\le C_\eta(n)
 \max_{\substack{W\ne\{0\}\\ W\text{ a lattice subspace of }\Lcal^*}}
 \det(\Lcal^*\cap W)^{-1/\dim W},
\end{equation}
and conjectures
\[
 C_\eta(n)\le \operatorname{polylog} n.
\]
The source immediately notes the homogeneous normalization: it suffices to
prove $\eta(\Lcal)\le\operatorname{polylog}n$ whenever every sublattice of
$\Lcal^*$ has determinant at least one.

Here $\eta(\Lcal)$ is the $1/2$-smoothing parameter, characterized by
\[
 \sum_{y\in\Lcal^*} e^{-\pi s^2\lVert y\rVert^2}=\frac32
 \quad\text{at }s=\eta(\Lcal).
\]

\section*{Later theorem}

The abstract and introduction of Regev--Stephens-Davidowitz explicitly say
that their main result proves Dadush's conjecture and cite the source paper.
Their Theorem 1.2 (``Reverse Minkowski Theorem''), on PDF page 3, states that
if $\Gamma\subset\R^n$ is a lattice satisfying
\[
 \det(\Gamma')\ge1
 \quad\text{for every sublattice }\Gamma'\subseteq\Gamma,
\]
then
\begin{equation}
 \sum_{y\in\Gamma}e^{-\pi t^2\lVert y\rVert^2}\le\frac32,
 \qquad t=10(\log n+2).
\end{equation}

\section*{Exact implication}

Set $\Gamma=\Lcal^*$.  Under the source's normalized hypothesis, the later
theorem is
exactly the smoothing inequality that yields
\[
 \eta(\Lcal)\le 10(\log n+2).
\]
The later paper also records the homogeneous form explicitly.  If
\[
 \eta_{\det}(\Gamma)
 =\max_{\Gamma'\subseteq\Gamma}
   \det(\Gamma')^{-1/\operatorname{rank}(\Gamma')},
\]
then its equation (1.3) gives
\[
 \eta^*(\Gamma)
 \le 10(\log n+2)\eta_{\det}(\Gamma),
\]
where $\eta^*(\Lcal^*)=\eta(\Lcal)$ in the source convention.  Comparing
this with the displayed source inequality proves the following.

\begin{theorem}
The Main Conjecture of Dadush--Regev is answered affirmatively in the later
literature, with the explicit estimate
\[
 C_\eta(n)\le10(\log n+2).
\]
\end{theorem}

This is stronger than the conjectured unspecified polylogarithmic bound.

\section*{Scope and verification}

This packet records a known answer rather than a new mathematical result.
It resolves the source's central Main Conjecture; it does not assert that
every stronger variant or related question discussed in the source is
settled.

The stored source PDF has 65 pages and the supporting PDF has 41 pages.
The source statement was checked in Conjecture 3.1, PDF page 14.  The answer
was checked in Theorem 1.2, PDF page 3, and in the later paper's equation
(1.3), which supplies the scaled formulation used above.  Cheap run-index
search found no prior packet for this exact identification.

\begin{thebibliography}{9}

\bibitem{DR16}
D. Dadush and O. Regev,
\emph{Towards Strong Reverse Minkowski-type Inequalities for Lattices},
arXiv:1606.06913 (2016).

\bibitem{RSD16}
O. Regev and N. Stephens-Davidowitz,
\emph{A Reverse Minkowski Theorem},
arXiv:1611.05979; Ann. of Math. 185 (2017), 233--273.

\end{thebibliography}

\end{document}
